\chapter[Did Colonial History End?]{Did Colonial History End with Colonial Empires?}
\section{A Coin with a Different Face}
First, pick up a coin---more precisely, a pair.

The first is a British Indian rupee from the 1940s. King George VI's profile faces left, surrounded by Latin titles. Its metal came from India. Indian clerks received it as pay, Indian farmers paid it as tax, passing the king's face from hand to hand daily. Most probably never paused to ask why their money carried a man who had never set foot on their land.

The second is an Indian coin from 1950. The king has disappeared, replaced by Ashoka's lion capital: lions standing back to back on a lotus base, a stone carving left by an Indian ruler over two thousand years earlier. Value, size, and feel are nearly unchanged; it still buys rice at market. But a quiet coup on its face is complete: the metal remains, the face changes.

A coin reveals independence's most visible meaning: new flags, stamps, portraits on money, and guards outside the governor's residence. India and Pakistan became independent in 1947; Ghana in 1957; seventeen African countries in 1960 alone. Across mid-century Asia and Africa, this "change of face" repeated as colonial flags descended and new nations presented calling cards at the UN's door.

Our question concerns the half a coin cannot answer: once the face changes, is colonialism over? The mint's die changes, but mines and cocoa fields remain, cargo ships sail to the same ports. A country can remove a king from its money without easily removing his companies, language, or casually drawn borders from daily life. When empire ends, does imperial history end too?

\section{Midnight and Trains}
Enter South Asia in August 1947.

Late on August 14, New Delhi's Constituent Assembly building blazed with light, crowded with people unwilling to go home. Nearly two centuries of British rule---first an East India Company governing a subcontinent as a corporation, then direct government---ended that night. At midnight, independence leader and incoming prime minister Jawaharlal Nehru rose to deliver the speech later known worldwide as "Tryst with Destiny." Outside, hundreds of thousands blew horns, beat drums, scattered petals over strangers, and filled the streets with car horns. Witnesses recalled a city where "nobody slept."

That same week, several hundred kilometers northwest in Punjab, another scene unfolded.

A British lawyer named Radcliffe had been sent to draw a dividing line. Muslim-majority northwestern areas and eastern Bengal would form Pakistan; the rest would go to India. Never previously in India, equipped with old maps and rough population figures, he drew thousands of kilometers of border in weeks, then boarded a plane home. Publication told Punjabi villages they were on "the other side." Hindus and Sikhs moved east with families; Muslims west. One of history's largest migrations began, ultimately uprooting tens of millions.

Then came trains. Refugee services from Delhi to Lahore and Amritsar to Delhi were intercepted. People on platforms stood on tiptoe for relatives. Doors opened onto carriages silent because nobody inside remained alive. Such "ghost trains" were not isolated in the latter half of 1947. Revenge killings, arson, and abductions accompanying Partition took lives estimated from several hundred thousand to over a million, without an agreed total.

Hold those images together: the same week, one subcontinent, Delhi's midnight horns and Punjab's platform silence. People celebrating freedom and fleeing death sometimes lived on the same street.

\section{The Real Questions}
Together the images reveal our real questions---more than one.

First: why did empire leave? A long-familiar British account says, "We transferred power honorably," as though independence were a graduation certificate issued in London, stamped before a shared photograph. Indian nationalists tell a different story: decades of marches, imprisonment, and fasting won it back; nobody issued anything. Both accounts cannot be right---or neither is entirely right.

Second: why did independence bleed? Expelling a foreign ruler should pit "us" against "them." Yet 1947's rivers of blood ran within "us," among Hindus, Muslims, and Sikhs who had shared village wells. How could a stranger's line, drawn in weeks, turn neighbors into killers?

Third, the title's question: does independence carry away everything colonialism left? Borders, railways, English, exports, armies, textbooks---which inherited things remain usable, and which are shackles renamed?

Do not answer yet. First hear voices from that night and the years after it.

\section{Voices around Midnight}
\textbf{Those transferring power.} The last viceroy, Mountbatten, was a royal relative and admiral assigned to manage the separation. His style was speed: a transition originally planned for several more years was advanced to August 1947. To him and many London officials this was pragmatic asset restructuring, not defeat. Postwar Britain was broke; Indian military and administrative spending looked bottomless. Staying cost money and invited abuse. Better negotiate respectably and exchange "rule" for "friendship and trade." Churchill, prime minister only years earlier, differed sharply. In the darkest wartime days he had publicly said, in essence, that he had not become prime minister to preside over the British Empire's dissolution. One Britain contained both views: withdrawal as wisdom or betrayal.

\textbf{Those taking power.} Nehru's midnight "tryst with destiny" redeemed generations of struggle for freedom. But where was Gandhi, the movement's acknowledged soul? Not at Delhi's festivities. He was in a besieged house in riot-torn Calcutta, urging Hindus and Muslims to put down knives. On a day of national joy, he chose fasting, prayer, and the place likeliest to ignite. To Gandhi, freedom that expelled Britain but brought neighbors killing neighbors lacked its most essential part.

\textbf{The other party to Partition.} Jinnah and the Muslim League had their own account. In a united, overwhelmingly Hindu India, Muslims would remain a permanent minority; democratic elections would never give them their turn. Statehood was insurance bought by a minority. Was the premium high? Punjab's silent platforms were part of it. Refugees reaching Karachi and Delhi later rebuilt cooking fires in new countries' slums. Histories called them "the cost of Partition"; they called themselves people who had lost their homes.

\textbf{Another African conversation.} Ten years later, Ghana saw a similar confrontation. On March 6, 1957, Nkrumah proclaimed independence in Accra, making Ghana the first sub-Saharan African country freed from colonial rule. His much-quoted words meant that Ghana's independence remained incomplete unless connected to all Africa's liberation. In preceding years, colonial officials and European commentators had repeatedly said, "You are not ready."

Pause to give them their fullest reasoning---not to agree, but to understand before rebutting.

Their case was roughly this: modern government needs accountants, engineers, judges, doctors, and functioning civil servants. Colonial education had trained only low-level clerks; local elites were too few. Hasty independence might collapse public finances, provoke intertribal war, and leave ordinary people worse off. That would be irresponsibility, not liberation. They could point to cases such as Congo's rapid post-independence turmoil.

There was truth here: new nations lacked trained personnel, as their leaders admitted. But a fatal flaw remained. Colonial rule had created that unreadiness by monopolizing higher education and senior posts, then saying, "You lack qualified people, so we must continue governing." It resembled tying someone to a chair and refusing release because they could not walk. Nkrumah and others answered simply: readiness can only be tested in independence. Besides, who sets the standard, the bound person or the one holding the ropes?

\textbf{An easily forgotten party.} Algeria's confrontation included not just France and the National Liberation Front but over a million European settlers. Born and raised there, with vineyards and bakeries there, they called it home. After independence in 1962, most crossed to France within months, many first setting foot in their supposed "homeland." Empire's ledger never had just colonizer and colonized columns. It also held these people moved by history.

\section{A Bridge for Thought: A Three-Rung Ladder of Why}
"Why did empire leave?" invites a trap: one answer. Defeat, awakened conscience, international pressure? Historians often use a portable tool, sorting causes into three levels like a ladder.

\textbf{Bottom rung: structure.} Slowly changing backgrounds do not cause events on particular days, but determine what becomes increasingly possible. Decolonization's structures included two world wars exhausting Europe and turning Britain from creditor to debtor; a century of colonial extraction accumulating subterranean resentment; and education and newspapers carrying "nation" and "self-determination" into colonial cities.

\textbf{Middle rung: pressure.} People organize and act upon structures. Without them, conditions merely smolder. Congress led decades of Indian mass movements involving millions. In 1930 Gandhi walked hundreds of kilometers to the sea to make salt against Britain's tax monopoly, inspiring nationwide action. Indonesian youths and nationalists rapidly assembled government amid Japanese occupation's wreckage. Algeria's FLN fought nearly eight years of guerrilla war from 1954. Nkrumah's "Positive Action," strikes, boycotts, and noncooperation, forced negotiations in Ghana. Different methods included negotiation, mass action, armed struggle, and more often combinations.

\textbf{Top rung: triggers.} A specific opening makes events happen now rather than later. Japan's 1945 surrender created a Southeast Asian power vacuum that Sukarno used to declare Indonesian independence. The Royal Indian Navy mutiny of 1946 showed London its own governing machine loosening. Egypt nationalized the Suez Canal in 1956; Britain and France invaded to recover it, then withdrew humiliated under American and Soviet pressure. Old empires learned they could not even reclaim a canal without superpower consent.

\textbf{Weather across all three: the international environment.} The postwar UN wrote self-determination into its documents. New superpowers America and the Soviet Union, for their own reasons, disliked the old empires retaining colonies. Colonizers looked around and found no audience applause.

Next time someone asks why something happened, do not rush to choose among answers. Build the ladder: what background, who pushing, what spark, what weather? Many arguments are simply people speaking from different rungs.

\section{A Declaration Entering the UN Record}
Now read a short primary document.

First hear midnight's words. On August 14, 1947, Nehru delivered "Tryst with Destiny" to the Constituent Assembly. Its most famous line means, in a rendering of the Chinese paraphrase:

\begin{quote}
When the midnight hour strikes, while the world sleeps, India will awaken to life and freedom. (Nehru, "Tryst with Destiny," Constituent Assembly address, August 14, 1947; rendered from the Chinese paraphrase of the English.)
\end{quote}
It is genuinely moving. But after Punjab's trains, hear its other side: at midnight the subcontinent was not "asleep." Thousands fled for their lives. India woke in the speech; real India and Pakistan entered bloodshed together. This does not call Nehru hypocritical. It reminds you to read primary materials beside the full scene of their creation. Podium words and platform bodies both belong to history.

Now a more official text. In December 1960, the UN General Assembly adopted Resolution 1514, the Declaration on the Granting of Independence to Colonial Countries and Peoples. Its central claim was that alien subjugation, domination, and exploitation deny fundamental human rights, and colonialism in every form must end immediately and unconditionally. Almost all Asian and African states voted for it. Old colonial powers including Britain, France, and Portugal abstained: neither opposed nor approved, a stance itself bearing historical witness.

Its significance was not whom it "liberated": paper declarations free nobody by themselves. It marked a change of referee. When Europeans partitioned Africa at Berlin half a century earlier, the premise was "who gets which part?" By 1960, the world's highest deliberative chamber began from "colonialism itself is illegitimate." Rules did not fall from heaven; decades of struggle punched and kicked them into the text. Remember: documents do not merely record history; they are battlefields where forces contest it.

\section{One Lowered Flag, Three Explanations}
Return to why empire left. Compare three genuinely different explanations, reasons and weaknesses included. Separate facts, colonial systems collapsing across decades; explanations, our present debate; contemporary understanding, transferors saying "we gave" and independence movements "we reclaimed"; and today's judgments. We compare the second.

\textbf{First: the arithmetic changed---empire no longer paid.} Decolonization was primarily metropolitan cost accounting. Two world wars left Britain deeply indebted; garrisons and administration grew costlier while extraction's relative returns shrank. Free trade and multinational companies also allowed business without occupying land. Once empire moved from profit to loss, liquidation became rational. Suez's 1956 retreat illustrates not unwillingness to stay but inability to afford it. This sober account explains some nearly peaceful withdrawals. Its weakness is turning decades of colonial sacrifice into a London accountant's entry. Why did it become unprofitable precisely then? Costs soared because people struck, marched, and fought guerrilla wars daily.

\textbf{Second: power was taken---struggle broke the barrier.} Restore colonized peoples as protagonists. India's nonviolent campaigns made British rule dearer daily. Eight years of Algerian gunfire required rotating hundreds of thousands of French troops and tore French politics apart. Vietnamese artillery and trenches directly defeated France at Dien Bien Phu. Without such action, no empire would voluntarily put down its calculator. Timing is the strongest evidence: independence came most "suddenly" where struggle intensified. Its limit is the broader timing: why did decades of struggle yield results together in the mid-twentieth century? Equally strong nineteenth-century resistance was usually crushed by fleets and guns.

\textbf{Third: the weather changed---new international referees.} At the widest scale, war exhausted old empires and produced two superpowers disliking colonial systems. The UN charter enshrined self-determination; the 1960 declaration turned colonialism from "normal" to "illegal." Both superpowers sought friends and disliked being called imperial accomplices. Portugal persisted until its own 1974 revolution finally brought withdrawal from Africa. This fills the timing gap, but crediting everything to climate makes prisoners, jungle fighters, and negotiators background figures again, merely fortunate in their weather.

Each account grasps a rung: the third structure, the first changing metropolitan pressures, the second colonial pressures. Beware any "one cause is enough" version. History has not set a single-choice question here.

\section{Three Assignments for New States}
Raising the flag only begins the homework. New states generally faced three problems written directly into their colonial inheritance.

\textbf{First: borders.} Many African boundaries were drawn with rulers by late-nineteenth-century European diplomats at Berlin: rivers sliced, peoples divided among three states, old enemies put together. Should independence redraw them? Intuition says colonial mistakes should be corrected. Yet the Organization of African Unity, founded in 1963, decided the following year to recognize inherited borders. Its practical reason: open the door to ethnic redrawing and almost every African country could fragment into endless territorial war. Independent states themselves froze colonial lines. Their justice remains disputed, but the shortcut "overturn every colonial legacy" was rejected by those first living with it.

\textbf{Second: the economic umbilical cord.} Independence did not redirect mines and plantations. Ghana still exported chiefly cocoa, Senegal peanuts, Zambia copper. Economies designed to supply metropolitan factories' raw materials remained intact. Others set commodity prices and sold manufactured imports; power lay elsewhere between sale and purchase. Currency imposed a firmer gate. The CFA francs of Francophone West and Central Africa remained linked to France, with most foreign reserves long required to be held in the French Treasury. The flag was local, but Paris retained a purse key. Reforms began only in recent years, and argument continues.

\textbf{Third: what to build the state with.} Colonial rule left an administrative skeleton of bureaucracy, police, and military designed to control, not serve. New governments had to transform it or be transformed by it. Language was unavoidable: what should schools, courts, and documents use? The colonial language was convenient, ready, internationally connected, but placed a foreign-language wall between elites and masses. Local languages offered familiarity and equality but required choosing among dozens or hundreds; excluded groups would ask why. India argued for decades and retained English as an associate official language. Tanzania promoted Swahili as a national common language. Both paths worked for some and exacted costs.

\section[The Cold War in the Tropics]{Somebody Else's Chessboard: The Cold War in the Tropics}
Before new states finished those assignments, a door opened onto the Cold War.

America and the Soviet Union confronted one another globally, each needing friends, bases, ports, and votes---things independent states possessed. Local conflicts quickly acquired superpower supply lines; local grievances became bloc wars.

Congo is the clearest scene. Independent from Belgium on June 30, 1960, it heard Prime Minister Patrice Lumumba frankly describe colonial suffering at the ceremony, offending the visiting Belgian king. Days later came army mutiny and Belgian-backed secession in mineral-rich Katanga. Lumumba sought UN help, then Soviet aid. Amid turmoil he was dismissed, arrested, and taken to Katanga for murder in early 1961. Later investigations confirmed deep Belgian official involvement; Belgium formally apologized in the twenty-first century. Barely half a year separated independence from the first elected prime minister's death. Local opponents fired the bullets, but the network arming them, undermining him at the UN, and pulling hidden strings crossed Brussels, Washington, and Moscow.

Vietnam followed another line. After France's 1954 defeat, the Geneva settlement temporarily divided Vietnam at the seventeenth parallel. America saw South Vietnam as the first domino preventing others falling and deepened involvement into a war exceeding a decade. Washington called it containing communism; Hanoi's supporters called it anticolonial struggle's second half. Both told coherent stories; Vietnamese people in rice fields bore most deaths.

Then Angola. After Portugal withdrew in 1975, three independence movements immediately fought. The Soviet Union and Cuba supported one side, America and South Africa another. Intermittent civil war lasted into the twenty-first century. Angolans contested their country's direction, but weapons, intelligence, and cash bore foreign emblems.

The lesson is twofold. Do not blame every new nation's disorder on colonial inheritance: the Cold War injected weapons and stakes that made negotiable conflicts intractable. Nor blame everything on the Cold War: superpowers often entered through existing colonial fractures of borders, ethnicity, and economy. Old wounds had not healed before salt was added.

\section[New Flags, Unchanged Shipping Routes]{A Deeper Challenge: New Flags, Unchanged Shipping Routes}
Now introduce a theory of the period after independence, developed by Latin American, Asian, and African economists: \textbf{dependency theory}.

Argentine economist Raúl Prebisch led with a simple question. Textbook free trade predicts shared enrichment through exchanging comparative advantages. Why, then, did the gap widen between poor raw-material exporters and rich manufacturers? His answer: the game began asymmetrically. The "center" sold machines, cars, medicines, technically sophisticated products with firm prices. The "periphery" sold cocoa, coffee, and copper, volatile commodities whose long-term purchasing power declined: the same shipload of cocoa bought fewer machines. The harder the periphery planted, the more firmly it remained peripheral. Immanuel Wallerstein later expanded this into world-systems theory: a global division of labor assigns some positions to earning profits and others to being profited from. Independence changes the flag, not the position.

Through this lens, puzzles become understandable. Why do children in leading cocoa producer Ghana rarely afford chocolate? Beans leave cheaply, become chocolate in Europe, and return dear; the richest value-adding stage happens elsewhere. Why must CFA countries consider outsiders before printing money to stimulate growth? Their currency anchor is not their own. Why do railways connect mines to ports rather than domestic cities? They were not built for domestic populations. Dependency theory locates colonialism's strongest part not in governors' residences but in routes, contracts, and price lists. None departed when flags fell.

Weigh its limits lest one extreme replace another. Blame all on external structures and local governments become forever innocent. Yet countries with similar beginnings took different paths; some East Asian economies moved from periphery toward center within decades, difficult for pure dependency theory to explain. Corruption, civil wars, and bad domestic decisions genuinely caused some poverty. Blaming "the world-system" alone is as lazy a single-cause explanation as saying poor countries are poor because people are lazy. Mature judgment holds both: the card table is crooked, but players' choices still count.

Add culture. French-descended scholar Edward Said argued in Orientalism that colonialism occupied not just land and resources but knowledge. It portrayed colonized peoples as lazy, mysterious, irrational, needing supervision. Newspapers, textbooks, and films repeated this until even the described peoples saw themselves through those lenses. If he is right, decolonization's hardest barrier lies neither in ports nor treasuries but in minds. A people raised on textbooks saying "your ancestors awaited our civilization" may need generations to remove the spectacles.

\section{Echoes: Trophies, Museums, and Your English Homework}
These old accounts are being settled around you now.

Enter European museums. The British Museum and several German collections display thousands of Benin Bronzes, exquisite sculptures looted when British forces conquered the kingdom of Benin, in present-day Nigeria, in 1897. For a century they were "treasures of human art"; increasingly they are called "stolen property." Germany formally transferred ownership of a group to Nigeria in 2022; others are returning objects while some delay. One bronze's place in a gallery literally poses our question: the robbery happened over a century ago, but while the object remains behind glass, the page is not turned.

Currency is another battlefield. Reform of the West African CFA franc sparks heated debate. Young people ask why their money remains linked to Paris after over sixty years of independence. Reformers and defenders both have reasons. A peg brings stability, vital to poor people; separation brings autonomy and its risks. This is not nostalgia but a live question.

Then language. In India, English opens good universities and jobs while remaining a contested colonial legacy. Young Algerians argue online over more Arabic or retained French. Language is never merely a tool. It is an unfinished battle continuing in classrooms and labor markets.

Finally, your desk. How many pages does your world history devote to colonialism, how many to decolonization? How many independence leaders can you name? Three chapters on building empires and three lines on ending them already express a position: conquerors matter, resisters are an epilogue. Said's spectacles sometimes lie between textbook pages.

\section{Your Question: What Counts as an Ending?}
Return to the coin.

In 1950, India removed the king and installed Ashoka's lions. Ceremony complete: empire legally ended. Yet the same country still uses English in higher education and a railway network shaped to carry goods overseas. Its hastily drawn border with Pakistan has brought several wars and remains heavily militarized.

What, then, marks colonial history's "end"?

The new flag? The last governor boarding ship? Local control of mines and banks? Looted objects returning? A generation no longer viewing itself through colonizers' spectacles? Or will it never "end," requiring successive generations to identify and clear its remains?

Before answering, weigh our voices again: Delhi's horns and Punjab's silence; Mountbatten's arithmetic and Gandhi's room; "not ready" and the person tied to a chair; Congo's prime minister dead within half a year; African states reluctantly freezing colonial borders; dependency theory's crooked card table and its not-entirely-blameless players.

There is no standard answer. But your answer shapes how you read international news, view museum objects, and approach a foreign language. Unfinished history need not mean permanent hatred. The bill lies open: neither those pretending it does not exist nor those forever clutching it have truly finished this chapter.
